\section{Frequency Analysis}

\subsection{Prinsip}

Dalam bahasa alami, setiap huruf memiliki frekuensi kemunculan yang khas. Untuk bahasa Inggris (contoh): E $\approx$ 12\%, T $\approx$ 9\%, A $\approx$ 8\%, O $\approx$ 7\%, I $\approx$ 7\%, N $\approx$ 7\%, S $\approx$ 6\%, ...

Untuk cipher substitusi monoalfabetik (termasuk Caesar), pemetaan huruf bersifat satu-satu. Huruf yang paling sering dalam ciphertext kemungkinan besar memetakan ke E, T, atau A. Dengan mencoba korespondensi dan memeriksa kemungkinan kata, penyerang dapat memecahkan cipher.

\subsection{Langkah Praktis}

\begin{enumerate}
  \item Hitung frekuensi setiap huruf dalam ciphertext
  \item Urutkan dari yang paling sering
  \item Coba pemetaan: huruf paling sering $\to$ E, kedua $\to$ T, dsb.
  \item Sesuaikan berdasarkan konteks (kata yang mungkin, bigram, trigram)
\end{enumerate}
